\chapter{Vitamin B1 (Thiamine) Deficiency --- Beriberi and Wernicke-Korsakoff Syndrome}
\index{Vitamin B1 (Thiamine) Deficiency --- Beriberi and Wernicke-Korsakoff Syndrome}
\label{sec:Vitamin B1 Deficiency}

\section{Overview}
Thiamine (vitamin B1) deficiency is a vitamin deficiency disorder affecting populations consuming polished white rice as a dietary staple, individuals with chronic alcoholism, hyperemesis gravidarum, after bariatric surgery, and in prolonged parenteral nutrition without thiamine supplementation. This genetic disorder results from specific gene mutations or chromosomal abnormalities. It may be present at birth or manifest later in life depending on the underlying mechanism. This metabolic disorder involves dysregulation of normal bodily processes, often requiring ongoing monitoring and management. It is increasingly common due to lifestyle and dietary factors. Genetic disorders result from abnormalities in an individual's DNA, ranging from single-gene mutations to chromosomal abnormalities. These conditions may be inherited from parents or arise from new mutations. Pattern of inheritance depends on whether the mutation is dominant, recessive, or X-linked. Genetic counseling helps families understand risks and make informed decisions.
\section{Causes}
This condition results from disruption of normal vitamin b1 (thiamine) deficiency --- beriberi and wernicke-korsakoff syndrome function.

\begin{itemize}[noitemsep]
  \item This condition happens because of thiamine deficiency
  \item Thiamine functions as a cofactor for several enzymes involved in carbohydrate metabolism, including transketolase, pyruvate dehydrogenase, and $\alpha$-ketoglutarate dehydrogenase.
\end{itemize}
\section{Risk Factors}


\begin{itemize}[noitemsep]
  \item Genetic predisposition and family history
  \item Advancing age
  \item Lifestyle factors (diet, exercise, smoking, alcohol)
  \item Environmental exposures
  \item Underlying medical conditions
  \item Medication use
\end{itemize}
\section{Signs and Symptoms}

\subsection{Common symptoms}
\begin{itemize}[noitemsep]
  \item You may experience two major syndromes: dry beriberi (symmetric peripheral neuropathy with stocking-glove sensory loss, muscle weakness, absent deep tendon reflexes, paresthesias), wet beriberi (high-output heart failure with dilated cardiomyopathy, peripheral vasodilation, swelling, tachycardia, warm extremities, bounding pulses, lactic acidosis), and Wernicke encephalopathy (acute neuropsychiatric emergency with confusion, loss of coordination, and ophthalmoplegia) which may progress to Korsakoff syndrome (irreversible amnestic disorder with anterograde and retrograde amnesia and confabulation).

  \item Pain or discomfort in the affected area
  \end{itemize}

\subsection{Emergency warning signs}
\begin{itemize}[noitemsep]
  \item Severe or worsening symptoms of vitamin b1 (thiamine) deficiency --- beriberi and wernicke-korsakoff syndrome require immediate medical evaluation
\end{itemize}
\section{Progression and Stages}
The condition may progress through different stages of severity over time. Early stages often involve mild or subtle symptoms that may be overlooked. Moderate stages involve more noticeable symptoms affecting daily function. Advanced stages may involve significant impairment and complications.
\section{Complications}
\begin{itemize}[noitemsep]
  \item Disease progression leading to worsening symptoms
  \item Secondary infections or injuries
  \item Side effects of treatment
  \item Reduced quality of life
  \item Emotional and psychological impact
\end{itemize}
\section{Diagnosis}
\begin{itemize}[noitemsep]
  \item Doctors usually diagnose this based on your symptoms and a physical exam
  \item Whole blood thiamine or erythrocyte transketolase activity may be measured but treatment should not be delayed.
  \item Comprehensive medical history and physical examination
  \item Laboratory tests to assess organ function
  \item Imaging studies to visualize affected structures
  \item Specialized testing depending on the affected system
  \item Consultation with relevant specialists
\end{itemize}
\section{Prevention}
\begin{itemize}[noitemsep]
  \item Maintain a healthy lifestyle with balanced diet and exercise
  \item Avoid known risk factors
  \item Attend regular medical check-ups for early detection
  \item Follow recommended screening guidelines
\end{itemize}
\section{Treatment Options}
\begin{itemize}[noitemsep]
  \item Treatment focuses on thiamine 100--500 mg IV/IM daily for 5--7 days followed by oral maintenance
  \item Wernicke encephalopathy requires immediate high-dose parenteral thiamine before glucose administration.
  \item Medications to manage symptoms and address underlying mechanisms
  \item Lifestyle modifications and self-care measures
  \item Physical therapy and rehabilitation
  \item Surgical intervention when indicated
  \item Regular monitoring and follow-up care
\end{itemize}
\section{Living with It}
\begin{itemize}[noitemsep]
  \item Follow your treatment plan as prescribed
  \item Maintain regular follow-up with your healthcare provider
  \item Make necessary lifestyle adjustments
  \item Seek support from family, friends, or support groups
  \item Stay informed about your condition
\end{itemize}
\section{Outlook}
\begin{itemize}[noitemsep]
  \item The outlook is generally positive for beriberi with treatment
  \item Wernicke encephalopathy is reversible with prompt treatment, but Korsakoff syndrome is irreversible.
\end{itemize}
